\chapter{Literature Review}

The rapid growth of digital technologies has significantly transformed modern healthcare systems. Hospitals are increasingly adopting digital platforms to improve operational efficiency, enhance patient engagement, and strengthen communication between patients and healthcare providers. Digital hospital systems, telemedicine platforms, and patient-centered health technologies play a major role in improving healthcare accessibility, quality of care, and overall patient satisfaction. This chapter reviews recent literature related to digital hospital systems, telemedicine, digital patient engagement, and patient-provider relationships.

The technologies used in digital hospitals include cloud-based platforms, integrated Electronic Medical Record (EMR) systems, mobile health applications, AI-based analytics, and secure hospital information systems. However, the study also identifies limitations such as high implementation cost, resistance from healthcare staff, technical complexity, data privacy concerns, and increased cognitive workload for clinicians. In some cases, excessive digital documentation reduced face-to-face interaction time with patients.

\section{Digital Hospital Systems and Patient Experience}

This study analyzes how fully digital hospitals influence both patient and clinician experiences. The functionality of digital hospital systems includes Electronic Health Records (EHR), digital medication management, clinical decision support systems, patient portals, and automated workflow systems. These systems aim to improve care coordination, reduce medical errors, enhance documentation accuracy, and enable real-time access to patient data. The study emphasizes improved communication, streamlined workflows, and enhanced patient safety as key outcomes of digital hospital implementation.
Technologies used in digital hospitals include cloud-based platforms, integrated Electronic Medical Record (EMR) systems, mobile health applications, AI-based analytics, and secure hospital information systems. Research emphasizes improved communication, streamlined workflows, and enhanced patient safety as key outcomes of digital hospital implementation. However, limitations include high implementation costs, resistance from healthcare staff, technical complexity, data privacy concerns, and increased cognitive workload for clinicians. In some cases, excessive digital documentation reduces face-to-face interaction time with patients.

\section{Telemedicine and Digital Patient Engagement}

This research focuses on redesigning telemedicine models to improve accessibility and patient engagement. The functionality of the proposed assisted telemedicine model includes remote consultations, virtual diagnosis support, digital appointment scheduling, remote patient monitoring, and digital follow-up care. The model particularly supports rural and underserved populations by integrating assisted telehealth centers where healthcare workers help patients access digital consultations.

The technologies used include video conferencing platforms, remote diagnostic devices, cloud-based health record storage, wearable health monitoring devices, and secure communication protocols. Despite its advantages, the system faces limitations such as internet connectivity issues, digital literacy barriers among patients, regulatory challenges, and limited physical examination capabilities. Additionally, data security and patient privacy remain major concerns.
\section{Digital Health Technologies in Patient Experience}

This scoping review examines various digital health technologies that influence patient experience. The functionality of these technologies includes mobile health (mHealth) apps, AI-driven chatbots, online appointment systems, digital feedback systems, and patient engagement portals. These tools aim to increase patient participation in treatment decisions, improve transparency in healthcare services, and provide personalized health recommendations.

The technologies discussed in the review include Artificial Intelligence (AI), Machine Learning (ML), wearable IoT devices, blockchain for secure data sharing, and mobile-based applications. However, the review highlights limitations such as technology accessibility gaps, interoperability issues between different hospital systems, ethical concerns in AI-based decision-making, and unequal access among elderly or low-income populations. Furthermore, long-term effectiveness data is still limited.
\section{Telemedicine and Patient-Provider Relationships}

This study explores how telemedicine impacts patient-provider relationships. The functionality of telemedicine platforms includes video consultations, secure messaging, digital prescription systems, and online medical record sharing. These platforms enhance convenience, reduce travel time, and allow continuous patient monitoring. The study highlights improved access to healthcare services and flexibility in appointment scheduling.

The technologies used include secure teleconferencing systems, encrypted communication channels, digital prescription software, and cloud-based data management systems. However, limitations include reduced physical interaction, potential miscommunication, technical glitches during consultations, and challenges in building trust through virtual platforms. In some cases, lack of emotional connection and non-verbal cues may affect patient satisfaction.
\section{Summary}

The proposed Bond-Health Hospital–Patient Interface System aims to overcome these limitations by providing:

\begin{itemize}
\item A centralized digital platform connecting hospitals and patients
\item Simplified appointment and specialist management
\item Secure patient data handling
\item A user-friendly interface for both hospital administrators and patients
\item Improved communication between healthcare providers and patients
\end{itemize}

Unlike many complex digital hospital systems, Bond-Health focuses on simplicity, accessibility, and integrated management of hospital services such as doctor registration, specialty management, laboratory, and pharmacy modules. By combining essential digital functionalities with an easy-to-use interface, the system enhances patient experience while maintaining operational efficiency.

Thus, this project contributes to the growing field of digital healthcare by offering a practical, scalable, and patient-centered hospital interface solution.